% Paper Summary: Let's Verify Step by Step (Lightman et al., 2023)

\begin{papersummary}{2023}{Let's Verify Step by Step}{Hunter Lightman, Vineet Kosaraju, Yura Burda, Harri Edwards, Bowen Baker, Teddy Lee, Jan Leike, John Schulman, Ilya Sutskever, Karl Cobbe}{Lightman et al. showed that rewarding each step of a chain of thought, rather than only its final answer, produces substantially more reliable mathematical reasoning than outcome supervision, and released PRM800K, 800,000 step-level human labels for training such reward models.}

\summaryheading{Key Ideas}
An outcome-supervised reward model (ORM) learns only whether a completion's final answer was correct; a process-supervised reward model (PRM) learns whether each intermediate step is correct, using per-step human labels. Training a PRM on PRM800K and using it to rerank many sampled solutions solved 78\% of a representative MATH test subset, outperforming an ORM of matched size and data. Process supervision has an advantage beyond raw accuracy: because every step is labeled, PRM-selected solutions are far less likely to reach a correct final answer through invalid reasoning, and active learning---using the current PRM to select which solutions are most valuable to label next---made the human annotation budget go substantially further.

\summaryheading{Follow-on Works}
Process supervision at scale required avoiding the cost of exhaustive human step-labeling; later work replaced human labels with automatically inferred ones (Monte Carlo estimation of step correctness, or a trained model that judges intermediate steps), making PRM-style training tractable without a labeling effort of PRM800K's size. \hyperref[paper:shao-2024-grpo]{Shao et al.'s GRPO (2024)} sidestepped the reward-model question differently, replacing both value and process reward models with the reward signal implicit in a group of sampled completions. \hyperref[paper:deepseek-2025-r1]{DeepSeek-R1 (DeepSeek-AI, 2025)} showed a full reasoning model trained this way with no process supervision at all, matching the reported performance of OpenAI's o1---whose own training methodology, unlike this paper's, was never disclosed.

\summaryheading{Lasting Contributions}
This paper established that reasoning quality can be trained, not just elicited by prompting: a model learns to avoid invalid steps because it was explicitly rewarded for avoiding them, not because a prompt asked it to think carefully. The result is widely read as the closest published account of the process-supervision ideas behind OpenAI's o-series reasoning models, whose full training methodology remains undisclosed. Reward-modeling reasoning credit at the step level, rather than only at the outcome level, became a standard design axis for every reasoning-training pipeline that followed, whether the per-step signal came from human labels, automated estimation, or an RL objective that made explicit process rewards unnecessary.

\end{papersummary}
